Confidence is not uniformly reliable. The ratio has lower Brier score than PCA in all synthetic configurations and SKAB, but a much worse Intel Brier score. Intel gives 0.068 with block interval [0.054, 0.082], compared with PCA's 0.017. Its smallest attainable window p-value is 0.059, so neither a 1\% nor a 5\% alarm level is representable. At level 0.10, 79 of 82 unchanged Intel cases alarm. The ratio's relative AP gain therefore does not establish a usable Intel alarm system. Normal cases include injection attempts that changed no observed coordinate.

The flow policy accepts useful synthetic repairs with no observed failures, but generally has lower coverage than PPCA. In S64N it accepts 125 of 320 cases, versus no PCA recommendation and 129 PPCA recommendations. Zero observed failures do not imply zero population risk, especially with dependent blocks. Intel accepts no repair and has undefined recommendation risk.

On SKAB the ratio accepts 123 repairs and makes 24 failed recommendations, all harmful. Actual risk is 0.195, with a wide block interval [0.000, 0.400], exceeding the calibration target in point estimate. PCA and PPCA risks are 0.371 and 0.183. H3 is not established as a general controlled-risk coverage advantage. Table~\ref{tab:repair} further shows that matched PPCA has lower repair CRPS and energy score in every configuration. Flow interval widths do not buy reliable real-data coverage. Figure~\ref{fig:repair} retains both successful and failed cases. Probability calibration concerns the designed injection prevalence, not operational physical-failure probabilities.
